\section{Conclusion}
\label{sec:conclusion}

When the target of surveillance is typical transmission or next period's value, the precision of a case-count estimate of the reproduction number grows with case volume at a diminishing rate. Beyond the balance scale, variation shared within a period dominates and each further doubling of cases removes at most a quarter of the remaining variance. In US influenza hospital admissions this scale was about 37 admissions per state-week on the scale of a renewal-type ratio. Conditional on a 3.6-day generation interval and on independent weekly shocks, the implied environmental variance was about 0.04 after calibration by simulation, and a large part of the shared variation was common to all states in a week. Surveillance systems that already report more cases per period than the balance scale gain more precision from more periods and from more regions with independent variation than from more cases per period. The evidence on reporting is narrower: under mandatory reporting, completeness contributed little to the intercept, and in an observational comparison high facility participation was associated with less volatile growth statistics. Empirical intercepts of apparent statistics should be reported together with the calibration that links them to transmission.
